\section{Tinjauan Pustaka}\label{sec:related}
\subsection{Root Cause Localization pada Sistem Microservice}
RCAEval adalah tolok ukur terbuka berisi 735 kasus kegagalan dari tiga sistem microservice, dilengkapi sejumlah baseline yang dapat direproduksi \cite{pham2024rcaeval}. Pendekatan berbasis struktur kausal yang lebih awal membangun graf dengan algoritma PC lalu merangking kandidat pada graf tersebut. RUN berargumen bahwa pendekatan itu mengabaikan urutan waktu dan menggantinya dengan neural Granger causal discovery \cite{lin2024run}. RUN terdiri atas tiga tahap: pre-training kontrastif backbone DLinear dengan dua potongan acak yang saling tumpang tindih dari deret yang sama sebagai pasangan positif tanpa pasangan negatif; jaringan peramalan per kasus dengan matriks atensi kausal yang dapat dipelajari, lalu diberi ambang dan dipangkas menjadi graf; serta PageRank dengan vektor personalisasi untuk merangking akar masalah dengan KPI pemicu (trigger) tertentu. RUN dievaluasi dengan HR@$k$ dan MRR pada data sintetis dan test bed Sock Shop (13 microservice).

ART adalah framework tak terawasi terpadu untuk deteksi anomali, triase kegagalan, dan RCL \cite{sun2024art}. ART menserialisasi metrik, log, dan trace menjadi jendela, memodelkan dependensi channel, temporal, dan call dengan Transformer encoder, GRU, dan GraphSAGE, dan melatih modelnya untuk memprediksi langkah waktu berikutnya pada data normal dengan loss galat kuadrat rata-rata. Saat diagnosis, selisih antara prediksi dan observasi membentuk representasi kegagalan tingkat instance dan tingkat sistem. Untuk RCL, instance dirangking berdasarkan kemiripan kosinus antara deviasi instance dan deviasi sistem, yang bertumpu pada pengamatan empiris bahwa instance akar masalah menyimpang lebih mirip dengan sistem secara keseluruhan. ART dievaluasi pada dua dataset, yaitu D1 (sistem e-commerce simulasi dengan 46 instance dan 210 kegagalan) dan D2 (sistem bank dengan 18 instance dan 133 kegagalan, tidak publik). Pembagian data dilakukan berdasarkan waktu (40\% kegagalan terakhir untuk pengujian). Metrik yang dilaporkan adalah Top1, Top3, dan AVG@5 dengan
$\mathrm{Top}K=\frac{1}{N}\sum_{i=1}^{N}\mathbb{1}[g_i\in P_{i,[1:K]}]$ dan $\mathrm{AVG@5}=\frac{1}{5}\sum_{K=1}^{5}\mathrm{Top}K$, dirata-ratakan dari lima kali pengulangan. Penulisnya mencatat dua keterbatasan yang relevan di sini: model terikat pada himpunan instance sehingga harus dilatih ulang ketika service berubah, dan kedua dataset berukuran kecil.

\subsection{Representasi Deret Waktu Self-Supervised}
TS2Vec mempelajari representasi pada tingkat semantik sembarang. Dua tampilan deret yang dipotong dan dimasking secara acak dikodekan oleh proyeksi masukan yang diikuti CNN dilated, dan pasangan positifnya adalah representasi timestamp yang sama pada kedua tampilan (contextual consistency). Loss menggabungkan kontras temporal dan kontras antar-instance, lalu diterapkan pada beberapa skala yang diperoleh melalui max pooling \cite{yue2022ts2vec}. TS2Vec dievaluasi pada 125 dataset UCR dan 29 dataset UEA untuk klasifikasi, dengan SVM dilatih di atas representasi yang dibekukan. TS2Vec juga dievaluasi pada peramalan dan deteksi anomali, dengan satu pengaturan memakai encoder yang dilatih pada FordA dan diuji pada Yahoo dan KPI. TimeMAE memotong deret menjadi sub-series yang tidak tumpang tindih dan memasking sebagian besarnya (60\% secara default). Transformer hanya mengodekan sub-series yang terlihat, sedangkan decoupled encoder berbasis cross-attention yang terpisah merepresentasikan posisi yang di-mask, sehingga tidak ada token mask yang masuk ke encoder utama. Tugas pretext-nya adalah masked codeword classification dengan tokenizer dan codebook yang dipelajari, serta masked representation regression terhadap target encoder yang diperbarui dengan momentum \cite{cheng2026timemae}. TimeMAE dievaluasi pada lima dataset klasifikasi dengan dua protokol yang disebutnya FineLast (encoder dibekukan) dan FineAll (fine-tuning penuh), dengan variasi proporsi label latih, serta pengaturan one-to-many ketika encoder yang di-pre-train pada HAR di-fine-tune pada empat dataset lainnya. Tidak satu pun dari kedua paper ini mengevaluasi lokalisasi di antara himpunan kandidat.

\subsection{Posisi dan Celah Penelitian}
Tabel~\ref{tab:related} membandingkan keempat metode dengan penelitian ini. Tiga perbedaan penting bagi penelitian ini.
\begin{enumerate}
\item \emph{Bentuk tugas.} TS2Vec dan TimeMAE menghasilkan satu representasi untuk satu deret, sedangkan RCL mengharuskan merepresentasikan setiap service dalam satu kasus lalu memilih satu di antaranya. Karena itu kami menambahkan scoring head di atas representasi per service (Bagian~\ref{sec:method}), yaitu pilihan rancangan yang tidak diambil dari kedua paper tersebut.
\item \emph{Peran label.} RUN dan ART tidak membutuhkan label, tetapi keduanya terikat pada dataset (RUN per kasus, ART per sistem) dan tidak menawarkan komponen yang dapat ditransfer dan sudah diuji. TS2Vec dan TimeMAE menghasilkan encoder yang dapat dipakai ulang, tetapi membutuhkan label untuk keputusan hilir. Studi transfer karena itu harus membandingkan kedua keluarga dengan aturan yang tidak merugikan salah satunya.
\item \emph{Bukti transfer.} TimeMAE menguji transfer antar-dataset klasifikasi dan TS2Vec antar-dataset deteksi anomali. ART dan RUN tidak menguji transfer, dan TS2Vec maupun TimeMAE belum pernah dievaluasi pada data RCA microservice. Kami mengadopsi istilah FineLast/FineAll serta rancangan fraksi label dari TimeMAE, dan menambahkan encoder yang diinisialisasi acak sebagai kontrol, analog dengan baseline berinisialisasi acak pada paper tersebut.
\end{enumerate}
Heuristik deviasi sederhana yang dipakai sebagai baseline di sini berkaitan dengan pengamatan ART bahwa deviasi anomali membawa informasi tentang akar masalah. Berbeda dari ART, heuristik kami memakai skor tetap dari robust $z$-score dan tidak dipelajari. ART tidak diikutsertakan dalam eksperimen. ART berfungsi sebagai acuan untuk definisi tugas, metrik (Top-$K$, AVG@5), dan pembahasan tentang pelatihan ulang ketika service berubah.
\begin{table}
\caption{Perbandingan metode terkait dengan penelitian ini (informasi diambil dari paper yang dikutip).}\label{tab:related}
\centering
\scriptsize
\setlength{\tabcolsep}{2.5pt}
\begin{tabular}{p{1.5cm}p{2.0cm}p{2.0cm}p{2.9cm}p{2.0cm}}
\toprule
 & Sinyal pembelajaran & Label dibutuhkan & Data yang dievaluasi & Transfer lintas dataset\\
\midrule
TS2Vec \cite{yue2022ts2vec} & Kontrastif hierarkis & Hanya untuk tugas hilir & UCR, UEA; peramalan; deteksi anomali & Terbatas (encoder dilatih pada FordA, diuji pada deteksi anomali Yahoo dan KPI), bukan RCA\\
TimeMAE \cite{cheng2026timemae} & Masked codeword dan regresi & Hanya untuk tugas hilir & 5 dataset klasifikasi & Ya (HAR $\rightarrow$ empat dataset), bukan RCA\\
RUN \cite{lin2024run} & Backbone kontrastif, graf Granger per kasus & Tidak ada & Sintetis; Sock Shop & Tidak diteliti\\
ART \cite{sun2024art} & Prediksi langkah berikutnya pada data normal & Tidak ada untuk RCL & D1 (publik), D2 (privat) & Tidak diteliti\\
Penelitian ini & TS2Vec, TimeMAE, RUN-encoder dengan scoring head & Label untuk fine-tuning & RCAEval RE1--RE3 & Ya (RE1, RE2, RE3)\\
\bottomrule
\end{tabular}
\end{table}

